\subsection{Datasets}

The AUDA analytical database contains over 28{,}000 human-reviewed extracted records across 195 countries and regions. There are 39 different groups of features, including demographic features such as population and GDP, waste management features such as plastic waste generation and mismanaged plastic waste, and plastic consumption features such as total resin consumption and polymer polyethylene terephthalate (PET) consumption.

Although the records are human-reviewed, the data points are still subject to measurement errors, reporting inconsistencies across countries and regions, and the use of different statistical methods by different organizations. Additionally, many data groups do not have enough data points to be used for modeling. Therefore, we select a subset of datasets that are suitable for modeling and analysis in this study.

\begin{table}[!ht]
	\caption{
		\centering
		Datasets selected for the correlation, importance, imputation, temporal forecasting, and pooled prediction experiments. Rows without a target variable are unlabeled feature sets.
	}
	\begin{tabular}{p{4cm} p{5cm} c c}
		\toprule
		\textbf{Dataset}         & \textbf{Target variable}   & \textbf{Years} & \textbf{Samples} \\
		\midrule
		PWDriverFeatureSet       & -                          & -              & 79               \\
		PWGPredictors            & Plastic Waste Generation   & -              & 695              \\
		YearTRC (Japan)          & Total Resin Consumption    & 2009-2017      & 9                \\
		YearTRC (United States)  & Total Resin Consumption    & 2009-2017      & 9                \\
		GlobalPlasticsProduction & Global Plastics Production & 1960-2019      & 59               \\
		YearPWG (Japan)          & Plastic Waste Generation   & 1995-2020      & 26               \\
		% YearPWG (US)             & Plastic Waste Generation   & 1990-2020      & 9                \\
		PWDrivers                & Plastic Waste Generation   & 2010-2019      & 617              \\
		%		PWDrivers (Japan)        & Plastic Waste Generation   & 2010-2019      & 10               \\
		%		PWDrivers (Slovenia)     & Plastic Waste Generation   & 2010-2018      & 6                \\
		\bottomrule
	\end{tabular}
	\label{tab:datasets}
\end{table}


This study primarily uses datasets in the AUDA analytical database, as summarized in Table~\ref{tab:datasets}. The \texttt{PWDriverFeatureSet} dataset is an unlabeled dataset comprising population, GDP, urban population, plastic waste generation, mismanaged plastic waste, and the proportion of plastic waste that is properly managed. The \texttt{PWGPredictors} dataset includes year, GDP, population, urban population, rural population, urban population ratio, and the rural population ratio as predictors for plastic waste generation. The \texttt{YearTRC (Japan)} and \texttt{YearTRC (United States)} datasets represent annual trends in total resin consumption for Japan and the United States, respectively. The \texttt{GlobalPlasticsProduction} dataset is obtained from \textit{Our World in Data} and provides the annual trend in global plastics production~\cite{owid_global_plastics_production}; the 1974 observation is absent from the original series. The \texttt{YearPWG (Japan)} dataset captures the annual trend in plastic waste generation in Japan. The \texttt{PWDrivers} dataset contains year, GDP, urban population, and rural population as predictors for plastic waste generation.
